
%% ============================================================================
%% PRE SUBMISSION CHECKLIST. Resolve every item before posting anywhere.
%%
%% 1. CVE IDENTIFIERS. Verify CVE-2026-42767, CVE-2026-45445, CVE-2026-14456
%%    and CVE-2023-0286 each resolve on nvd.nist.gov and that the affected
%%    component and version range match what the manifest asserts. A CVE ID in
%%    a published paper that does not resolve is unrecoverable.
%%
%% 2. SECONDARY CITATIONS. Cofano, Benedetti and Dann are currently cited as
%%    characterised by Churakova et al. Either obtain and read the primary
%%    sources and cite them directly with full bibliographic detail, or keep
%%    the present "as surveyed by" framing. Do not fill in a full citation for
%%    a paper you have not read.
%%
%% 3. AI ASSISTANCE DISCLOSURE. ACM policy requires disclosing use of
%%    generative AI in preparing the work. Add a statement in the
%%    acknowledgements describing what was used and for what, and confirming
%%    the author verified all technical content. Check the target venue's
%%    current wording before submitting.
%%
%% 4. AFFILIATION. Currently "Independent". If this work was done as a student
%%    at an institution, check that institution's IP and publication policy
%%    before claiming or omitting the affiliation.
%%
%% 5. VENDOR NAMING. Section 6 deliberately does not name the scanner or SBOM
%%    generator, pending the author's own 14 day private notice and right of
%%    reply policy. Once that period has elapsed with the reports public,
%%    decide explicitly whether to name them. Do not leave this by default.
%%
%% 6. DOUBLE BLIND. Remove the artifact footnote and the author block for any
%%    anonymous submission. The repository URL contains the author username.
%%
%% 7. DOI. Insert the archived dataset DOI in the artifact footnote once
%%    reserved.
%% ============================================================================
%% LADING CP-11
%% Submission wrapper, ACM sigconf.
%% For a double blind venue: add anonymous=true to the class options and
%% replace the author block with \author{Anonymous Author(s)}.
\documentclass[sigconf,nonacm]{acmart}

\usepackage{booktabs}
\usepackage{array}
\usepackage{url}
\usepackage{listings}

\lstset{
  basicstyle=\ttfamily\scriptsize,
  breaklines=true,
  frame=single,
  columns=fullflexible,
  keepspaces=true,
  showstringspaces=false
}

\begin{document}

\title{Binary Grounded VEX Clearance Fails Three Times: Upstream of Evidence, on
Internal Symbols, and on Version Applicability}
\subtitle{A pre registered negative result, and two corrections to it, across 55 scanned
artifacts and 15,641 scanner findings}

\author{Gautam Khosla}
\affiliation{%
  \institution{Independent}
  \country{}
}
\email{}

\begin{abstract}
Regulatory regimes including the EU Cyber Resilience Act make vulnerability
reachability load bearing: a flaw in a component that cannot be reached in a
shipped product does not start the mandatory reporting clock. This creates demand
for automated binary grounded VEX clearance, in which a tool inspects the
artifact that ships and decides whether the vulnerable code is present.

We built such a pipeline and measured where it fails. Across 55 scanned
directories covering 53 catalogued shipped artifacts, including 12 firmware
images, and 15,641 scanner reported CVEs, we attribute every decision to the
pipeline stage at which it terminated. Two kill tests, written into the
repository before any code existed and never adjusted, resolve against the
approach.

First, the pipeline terminates upstream of evidence. 91.93\% of attributed
findings stop at identity resolution and a further 7.67\% at manifest lookup,
before any binary is opened. Zero findings terminate at a symbol table stage,
despite 83\% of corpus binaries being stripped. The coverage kill test is
therefore recorded as not evaluable rather than failed: the instrument never
reached the question it was built to answer. Termination is not uniform across
packaging ecosystems, with RPM based artifacts failing earlier and more
completely than Debian based ones.

Second, where the pipeline did decide, its evidence could not support its claim.
All 35 clearances justified \texttt{vulnerable\_code\_not\_present} by absence of
the vulnerable function from the dynamic symbol table. Those functions are
internal to OpenSSL and never appear in any export map, so absence was guaranteed
a priori and the verdict was unfalsifiable rather than merely wrong. Hand
verification of 100 statements found 20 of 20 clearances false and 80 of 80
refusals correct.

Third, and found only after the second result was written up, the pipeline never
computes whether a CVE applies to the version on the artifact. Two of four
manifest CVEs do not apply to the versions where they were decided. This
reclassifies 7 of the 35 clearances as correct for a reason never computed,
leaving 28 unsound, and corrects 4 of the 20 hand labels, leaving 16. The
soundness test still fails: its bar was one. The guard we added in response to
the second result closed symbol observability and left version applicability
untouched.

We publish the corpus, the hand labelled ground truth, and all figures as re
derivable artifacts. The binding constraint on binary grounded VEX clearance is
CVE to component identity, not evidence availability, and symbol table absence is
sound evidence only for exported functions.
\end{abstract}

\keywords{SBOM, VEX, software composition analysis, vulnerability management,
firmware, Cyber Resilience Act, negative result, reproducibility}

\maketitle

%% ARTIFACT LINK. Present in the camera ready and on arXiv.
%% REMOVE FOR ANY DOUBLE BLIND SUBMISSION: the URL contains the author's
%% username and would deanonymize the paper. SCORED permits artifact links
%% but requires that neither the link nor the site reveal the authors.
\renewcommand{\thefootnote}{}
\footnotetext{Code, corpus catalogue with hashes, all decision records, evidence
bundles, the hand labelled ground truth and the figure provenance audit:
\url{https://github.com/GautamTalksDev/Lading}. Every figure in this paper
regenerates from \texttt{scripts/rederive-results.sh}. Archived dataset DOI to
follow.}
\renewcommand{\thefootnote}{\arabic{footnote}}


\section{Introduction}

A vulnerability scanner pointed at a shipped firmware image returns hundreds of
CVE rows. Most are irrelevant to the binary that actually ships: the affected
code was compiled out, never linked, or belongs to a subsystem the product does
not use. Under the EU Cyber Resilience Act, and under the customer security
questionnaires that already function as de facto regulation, somebody must
justify each row. The Vulnerability Exploitability eXchange (VEX) formats carry
those justifications in machine readable form, and one of the permitted
justifications, \vcnp{}, is a claim about the
artifact rather than about its dependency graph.

That claim invites automation. If the vulnerable function is named and the
shipped binary is on disk, the question looks mechanical: is the function there
or not. This paper reports what happened when we built that tool and measured it
against thresholds fixed in advance.

We built \textsc{Lading}, a refusal first pipeline that emits a VEX statement
only when it can point at specific bytes on disk and re derive the reasoning
offline. Anything short of that produces \texttt{under\_investigation}. There is
no scoring, no heuristic, and no model in the decision path. Two kill tests were
written into the repository README before any application code existed, and
neither threshold moved at any point across three instrument passes.

Both resolved against the tool.

\subsection{Contributions}

\textbf{A stage attributed termination map.} We instrument every refusal with the
pipeline stage that produced it and report the histogram over 15,385 attributed
findings on 55 scanned directories covering 53 catalogued shipped artifacts.
99.60\% of findings terminate at identity resolution or manifest lookup. Zero
terminate at a symbol table stage. The literature establishes that identity and
naming are hard; we locate where that hardness stops a binary grounded pipeline,
and show it stops upstream of the evidence question entirely.

\textbf{A named, reproducible soundness failure for symbol based clearance.}
Symbol table absence is a natural
justification for the claim that vulnerable code is not present. It is unsound
for internal functions in a way that is invisible from outside: the check returns absent on every build
regardless of whether the code ships. We found this by hand verifying our own
tool's output, and nothing in the tool was capable of detecting it. We give the
mechanism, the disassembly evidence, and a contrast case where the same check is
sound.

\textbf{A published, hand labelled ground truth.} One hundred statements sampled
from real pipeline output on real artifacts, each labelled by hand against the
binary. We are not aware of an equivalent labelled set published for binary
grounded VEX clearance.

We also report a secondary interoperability result encountered while building the
corpus (Section~\ref{sec:interop}), in which one scanner reading one image
through two input paths disagrees in both directions with no error emitted.

\subsection{What this paper does not claim}

It is not a claim that scanners disagree, which is established. It is not a claim
that SBOM naming loses vulnerabilities, which is quantified elsewhere. It is not
a corpus scale demonstration that symbol evidence is insufficient, because we
never ran symbol rules at corpus scale. Section~\ref{sec:related} states these
non claims explicitly and Section~\ref{sec:limits} bounds what we did measure.

\section{Background}

\subsection{Reachability as a regulatory object}

The Cyber Resilience Act obliges manufacturers to document vulnerabilities in
products placed on the EU market and to report actively exploited vulnerabilities
on a fixed clock. A vulnerability present in a bundled component but not
reachable in the product does not carry the same obligation. The consequence is
economic: for a firmware vendor whose scanner emits hundreds of findings per
release, the difference between justifying each by hand and justifying them
automatically is a line item.

Among the permitted \texttt{not\_affected} justifications, most are claims about
program behaviour: that the code cannot be reached in execution, that an
adversary cannot control it, that mitigations exist. Those are hard to establish
and harder to audit. One justification is different in kind:
\vcnp{} asserts a fact about the artifact, and is
the only one a deterministic inspection of the shipped bytes could hope to
settle.

\subsection{The design we tested}

\textsc{Lading} takes a scanner finding and walks a chain. \emph{Identity
resolution} maps the scanner's package reference to an upstream component; a
Debian scanner reports \texttt{libssl3} while the advisory is filed against
\texttt{openssl}. The match is graded from \texttt{none} to \texttt{exact} and
only sufficient matches proceed. \emph{Manifest lookup} consults a curated
manifest tying the CVE to the upstream functions the vulnerability lives in, each
with a reviewed provenance URL for the fix commit. \emph{Symbol evidence} locates
the shipped binaries and determines whether the named symbols are present.
Finally the pipeline \emph{decides or refuses}, emitting
\texttt{not\_affected} with a content addressed evidence bundle, emitting
\texttt{affected}, or refusing with a reason code.

Two rules can clear a finding. \textbf{D01} (\texttt{component\_not\_present})
fires when the SBOM names a component whose identity symbols and library
dependencies are absent from the artifact. \textbf{D02}
(\vcnp{}) fires when the component is present but
the manifest named vulnerable symbol is not. A third rule, \textbf{D04}, emits
\texttt{affected} when a vulnerable symbol is observed.

The design is deliberately refusal biased. Three VEX justifications are forbidden
in source and enforced by a CI check that greps for them: the tool will not claim
non reachability in an execute path, non controllability by an adversary, or the
presence of inline mitigations. None can be established from deterministic
inventory.

\subsection{Pre registration}

Two kill tests were written into the repository README at the first commit,
before any application logic existed.

\textbf{KT-1 (coverage).} Across at least 40 real shipped artifacts, if fewer than
30\% of scanner reported CVEs resolve to a decidable \texttt{not\_affected} with
a re derivable evidence bundle, the approach does not work at usable scale.

\textbf{KT-2 (soundness).} In 100 hand verified statements, even one false
\texttt{not\_affected} means the tool is unsound. Not a rate. One.

The asymmetry is deliberate. A false refusal costs a reviewer time. A false
clearance removes a real vulnerability from somebody's radar and hands them a
document asserting it was checked.

A third clause governs how KT-1 may be read. It matters enough to quote in full,
because Section~\ref{sec:kt1} turns on it.

\begin{quote}
\textbf{\S11, Kill test evaluability.} KT-1 tests whether $\geq$30\% of scanner
reported CVEs resolve to a decidable \texttt{not\_affected} with a re derivable
evidence bundle. If refusal stage attribution shows the instrument did not reach
symbol table or evidence evaluation, that is zero symbol table refusals and
$\geq$99\% of findings terminating at identity resolution or manifest lookup,
KT-1 is \textbf{not evaluable}, not fail; the measured decided rate must not be
read as an answer to the pre registered question. Three consecutive not evaluable
results on the same corpus after instrument passes intended to restore
evaluability constitute a pattern; the pre registered action is stop.
\end{quote}

This clause prevents a specific self deception: reporting a low decided rate as a
finding about symbols when the pipeline never evaluated a symbol. It also caps
how many times the instrument may be repaired before the result is the repair
record rather than the measurement.

\section{Related work}
\label{sec:related}

\paragraph{Cross tool VEX consistency.} Churakova et al.~\cite{churakova} measure
agreement among seven VEX generation tools across 48 Docker Hub images, comparing
reported vulnerability sets with Jaccard and Tversky overlap and running input
type controls for direct image scanning against native SBOM ingestion. They
attribute most divergence to differences in vulnerability database coverage,
reporting a Pearson correlation of 0.88 between database overlap and report
overlap. Their design measures disagreement between tools. It does not score
false \texttt{not\_affected} claims, does not inspect shipped binaries, and by
construction cannot isolate one tool disagreeing with itself across two ingestion
paths, which is the shape of the result in Section~\ref{sec:interop}.

\paragraph{VEX scope and component variants.} Rasheed et al.~\cite{rasheed} study
two mismatch patterns in SBOM based composition analysis on four hand built Java
and Maven fixtures. In their first study, Grype and CVE-Bin-Tool apply no
suppression relative to the fixture VEX while Trivy suppresses all four cases
including reachable ones. In their second, scanners largely ignore CycloneDX
\texttt{pedigree.variants} lineage metadata. \textbf{We cite that inert versus
overbroad split; we do not claim it.} Their work sits one layer above ours: they
measure whether a scanner honours a VEX statement's scope. We measure whether the
evidence behind such a statement can be produced at all, on native binaries
rather than Java, and where the attempt terminates.

\paragraph{Identity and naming as a known bottleneck.} That SBOM identity is
lossy is established. As surveyed by Churakova et al.~\cite{churakova}: Cofano et
al.~\cite{cofano} analyse the accuracy of SBOM generators in the Python
ecosystem; Benedetti et al.~\cite{benedetti} measure the downstream effect and
estimate that roughly 20\% of vulnerabilities are identifiable from generated
SBOMs at current generator precision; and Dann et al.~\cite{dann} report a low
true positive rate on a Java oriented scanner test suite. We report these figures as characterised in that survey rather than
as our own reading of the primary sources. Our contribution is not the discovery
that
identity is hard. It is the measurement of \emph{where} identity terminates a
pipeline that would demand binary proof if it ever arrived there, and the
observation that it terminates so far upstream that the binary question never
arises.

\paragraph{Non claims.} Explicitly, this paper does not claim: first report that
scanners or VEX tools disagree; discovery of the Grype inert and Trivy overbroad
split; discovery that SBOM naming loses vulnerabilities; corpus scale proof that
symbol evidence is insufficient; a positive clearance rate or a KT-1 pass or fail
verdict; general superiority of any scanner; or soundness results beyond OpenSSL
D02 clearances.

\section{Method}

\subsection{Corpus}

The catalogue (\texttt{ARTIFACTS.yaml}, version \texttt{cp11-3}) holds 53 product
artifacts, all present with a recorded SHA-256, source URL and fetch timestamp.
Artifacts that could not be obtained are recorded as absent rather than replaced
silently; where a stand in was used it carries an explicit
\texttt{substitute-container} class label so that no result attributes a
container measurement to firmware. Table~\ref{tab:corpus} gives the composition.

\begin{table}[t]
\caption{Catalogued corpus composition, 53 product artifacts. Scanning covers 55
result directories: these 53 plus two auxiliary single binary rescan directories
discussed in Section~\ref{sec:denominator}.}
\label{tab:corpus}
\centering
\begin{tabular}{lr}
\toprule
Class & Artifacts \\
\midrule
firmware & 12 \\
oci-app & 12 \\
oci-base & 9 \\
static-binary & 8 \\
substitute-container & 7 \\
rtos-sdk & 3 \\
openwrt & 2 \\
\midrule
Total & 53 \\
\bottomrule
\end{tabular}
\end{table}

The 12 firmware class artifacts were obtained by hand: nine OpenWrt release
images across four release series, two Yocto reference builds, and one Netgear
consumer router image, each with a recorded hash and fetch log. Several vendor
sources could not be retrieved at all, including a TP-Link GPL portal and Netgear
GPL source archives; those failures are recorded in the fetch log.

\paragraph{The corpus is not version uniform.} Every component version claim in
this paper is per artifact, established from the version string in the shipped
shared object and corroborated by package metadata. This needs stating because
an earlier draft of our own analysis treated the corpus as carrying a single
OpenSSL version, and it does not. Across the artifacts that produced a decision
on the four manifest CVEs we observe seven distinct versions: 1.0.2h on the
Netgear firmware, 3.0.11, 3.0.13, 3.0.19 and 3.0.20 across Debian and Ubuntu
based images, and 3.5.6 and 3.5.7 on newer Debian and Alpine images. One
artifact additionally ships a second, differently versioned copy of the library
under a separate prefix. Section~\ref{sec:version} reports what followed from
not having measured this earlier.

\subsection{Pipeline and tool versions}

Each artifact is extracted to a profile rootfs, scanned with Grype 0.115.0, and
the findings passed to \textsc{Lading}, which emits one row per finding to
\texttt{decisions.jsonl} together with an evidence bundle for decided rows.
Aggregation and re derivation run from a single script. Corpus numbers flow
through Grype only.

All tool binaries were verified against official release checksums. One note on
provenance hygiene: Trivy, used in the interoperability measurement of
Section~\ref{sec:interop} but not in the corpus scan, was compromised twice in
2026. The version used post dates that window and its hash matches the official
tarball.

\subsection{Two traversals of the corpus}
\label{sec:traversals}

The instrument computes corpus figures by two distinct traversals, and every
count in this paper is labelled with the one that produced it. The
\emph{catalogue traversal} walks \texttt{ARTIFACTS.yaml} and writes the per
artifact binary profile. The \emph{results traversal} walks the scan result
directories and writes the aggregate telemetry. They do not cover the same set,
and summing one will not reproduce the other. Section~\ref{sec:denominator}
reports the difference and its consequences rather than reconciling it, because
the two answer different questions and both are published.

\subsection{Refusal stages}

Every refusal carries a reason code and every reason code maps to exactly one
pipeline stage. The stage order in Table~\ref{tab:stages} is read from the
implementation rather than assumed, which matters: provenance (S4) is evaluated
after global symbol table usability (S3) and before the per component stripped
gates (S5), not before symbols as a naive reading would suggest.

\begin{table}[t]
\caption{Pipeline stages and their reason codes.}
\label{tab:stages}
\centering
\footnotesize
\setlength{\tabcolsep}{3pt}
\begin{tabular}{@{}l>{\raggedright\arraybackslash}p{1.65cm}>{\raggedright\arraybackslash}p{4.15cm}@{}}
\toprule
ID & Stage & Reason codes \\
\midrule
S1 & identity resolution & \texttt{purl\_match\_\-insufficient}, \texttt{mapping\_\-probable\_only}, \texttt{no\_identity\_\-mapping}, \texttt{version\_\-underivable} \\
S2 & manifest lookup & \texttt{manifest\_no\_entry}, \texttt{manifest\_\-probable\_only} \\
S3 & symbol table usability & \texttt{symbol\_table\_\-unusable} \\
S4 & provenance gate & \texttt{provenance\_\-unverified} \\
S5 & stripped gates & \texttt{stripped\_\-insufficient\_dynsym}, \texttt{stripped\_\-static\_binary} \\
S5b & symbol observability & \texttt{symbol\_not\_\-observable} \\
S6 & identity unverified & \texttt{identity\_\-unverified} \\
S7 & evidence evaluation & decided (empty code) or \texttt{default\_\-insufficient} \\
\bottomrule
\end{tabular}
\end{table}

Stage S5b did not exist when the corpus was first scanned. It was added in
response to the soundness finding of Section~\ref{sec:soundness}; its
introduction is dated and reported rather than folded silently into the
instrument.

\subsection{Ground truth and the frozen snapshot}
\label{sec:frozen}

The KT-2 evidence set holds 100 statements sampled with a fixed seed from real
\texttt{decisions.jsonl} output: 80 refusals and 20 clearances. Each was labelled
by hand against the binary, not by reading the tool's own evidence back to
itself.

One methodological point governs how every number in Section~\ref{sec:results}
should be read. Each row records the verdict the pipeline emitted \emph{before}
the S5b guard was added. KT-2 asks whether those clearances were sound when
emitted. Rescoring the same rows against the corrected engine would evaluate a
different instrument, one modified in direct response to the finding under test,
and would void the pre registration. The snapshot is therefore frozen, and post
correction figures are reported separately and labelled throughout. Where a
number could belong to either snapshot we state which.

\section{Results}
\label{sec:results}

\subsection{Result 1: termination upstream of evidence}

\begin{table}[t]
\caption{Termination stage for all 15,385 attributed findings in
\texttt{decisions.jsonl}, post guard. Results traversal.}
\label{tab:histogram}
\centering
\begin{tabular}{@{}llrrr@{}}
\toprule
ID & Stage & Count & \% & Cum. \% \\
\midrule
S1 & identity resolution & 14,144 & 91.93 & 91.93 \\
S2 & manifest lookup & 1,180 & 7.67 & \textbf{99.60} \\
S3 & symbol table usability & \textbf{0} & 0.00 & 99.60 \\
S4 & provenance gate & 21 & 0.14 & 99.74 \\
S5 & stripped gates & \textbf{0} & 0.00 & 99.74 \\
S5b & symbol observability & 35 & 0.23 & 99.97 \\
S6 & identity unverified & 0 & 0.00 & 99.97 \\
S7 & decided & 5 & 0.03 & 100.00 \\
\bottomrule
\end{tabular}
\end{table}

Table~\ref{tab:histogram} is the central measurement of this paper. Two entries
are zero, and those zeros carry the result.

Of 15,385 attributed findings, 15,324 (99.60\%) terminate at identity resolution
or manifest lookup. Sixty one reach symbol evaluation or later. Five reach a
decided verdict.

The dominant reason codes are \texttt{purl\_\-match\_\-insufficient} (6,591),
\texttt{mapping\_probable\_only} (4,183), \texttt{no\_identity\_mapping} (3,370)
and \texttt{manifest\_no\_entry} (1,180). The first three are identity failures
of different shapes: the package reference cannot be matched to a component with
sufficient confidence, the available alias is probable rather than reviewed, or
no alias exists at all.

\paragraph{The zeros were the expected wall.} Table~\ref{tab:profile} shows why
this surprised us. Across the corpus, 83.0\% of inventoried binaries are
stripped. We designed around that constraint, built the stripped gates at S5, and
expected them to dominate the histogram. They never fired once. The pipeline died
so far upstream that the condition we had engineered for was never tested.

\begin{table}[t]
\caption{Binary profile by stratum, catalogue traversal, summed from the
published per artifact profile. Containers are defined as oci-base plus oci-app
plus substitute-container, which is the grouping the profile tool uses. The
firmware stratum is materially less stripped and more static than the container
stratum, so it is where symbol evidence would have been most viable.}
\label{tab:profile}
\centering
\begin{tabular}{@{}lrrr@{}}
\toprule
Stratum & Binaries & Stripped & Static \\
\midrule
Firmware & 4,236 & 75.5\% & 42.9\% \\
Containers & 18,724 & 96.6\% & 2.7\% \\
\midrule
All catalogued & 27,658 & 83.0\% & 18.8\% \\
All scanned & 27,670 & 83.0\% & 18.8\% \\
\bottomrule
\end{tabular}
\end{table}

The final two rows of Table~\ref{tab:profile} are the two traversals of
Section~\ref{sec:traversals}. They differ by 12 binaries and both are published;
Section~\ref{sec:denominator} accounts for the difference exactly.

\paragraph{Termination is not uniform across ecosystems.} On the three RPM based
artifacts in the corpus (Fedora 40, Rocky Linux 9, UBI9 minimal),
\texttt{no\_identity\_mapping} alone accounts for 19, 318 and 61 refusals
respectively. This is not an RPM source package normalization failure; it is the
absence of any alias entry at all. The alias seed carries ten rows, nine of them
probable, and is Debian shaped. RPM artifacts fail at identity earlier and more
completely than Debian based ones.

The practical consequence is that identity resolution is not one problem solved
once. It appears to require separate work per packaging ecosystem, and a table
built against one distribution family will underperform on another without
failing loudly.

\paragraph{Firmware stratum.} On the 12 firmware artifacts (5,553 findings),
95.50\% terminate at S1 and 4.27\% at S2. One finding is decided. The stratum
with the most favourable binary profile for symbol evidence is also where
identity resolution performs worst, consistent with firmware vendors shipping
components whose provenance the scanner infers rather than reads.

\subsection{KT-1: not evaluable}
\label{sec:kt1}

\begin{table}[t]
\caption{Decided coverage over the 15,641 finding aggregate. Neither snapshot
answers the pre registered question.}
\label{tab:kt1}
\centering
\begin{tabular}{@{}lrr@{}}
\toprule
Metric & Post guard & Pre guard \\
\midrule
Decided & 5 & 40 \\
Decided rate & 0.032\% & 0.2557\% \\
\texttt{not\_affected} with bundle & 0 & 35 \\
\texttt{affected} & 5 & 5 \\
\bottomrule
\end{tabular}
\end{table}

The measured decided rate is 0.032\% post guard, with a 95\% Wilson interval of
0.0137\% to 0.0748\%. The pre guard snapshot is 0.2557\%.

Both are far below the 30\% bar, and it would be easy, and wrong, to record KT-1
as a fail. The pre registered evaluability rule forbids it. Zero findings
terminated at a symbol table stage and 99.60\% terminated before any symbol gate.
The pipeline did not test the hypothesis KT-1 encodes and then lose; it never ran
the test. A fail label would assert a measurement we did not make.

\textbf{KT-1 is recorded as not evaluable.} An earlier reading of the same data
labelled it fail on the 0.2557\% figure. That label is formally withdrawn, and
the withdrawn reading is preserved in the evidence trail rather than deleted,
because a corrected record that hides the correction is worth less than one that
shows it.

This is the third consecutive not evaluable result on this corpus, across three
instrument passes each intended to restore evaluability. Under the pre registered
rule, three is a pattern and the action is stop.

\subsection{The denominator, stated precisely}
\label{sec:denominator}

KT-1 is computed over 15,641 findings drawn from 55 scanned directories, while
the stage histogram covers 15,385 findings. The gap is 256 and it is fully
accounted for. Three directories carry findings but produced no attributed
decision rows: two auxiliary single binary rescan directories with 144 and 109
findings, and one container base image with 3. Those sum to 256 exactly. A
further seven directories have empty decision files and contribute nothing,
because they carry no findings at all.

The two rescan directories are not product artifacts. They are not in
\texttt{ARTIFACTS.yaml} and the aggregate assigns them to an \texttt{unknown}
class precisely because the lookup fails. Their 253 findings are nonetheless
inside the pre registered KT-1 denominator by construction of the aggregation
script. Excluding them gives 15,388, and the decided rate moves from 0.0320\% to
0.0325\%, which changes nothing about the verdict but should be on the record
rather than discovered by a reader.

The same accounting explains the 12 binary difference between the two traversals
in Table~\ref{tab:profile}. The results traversal includes those two rescan
directories (one binary each) and 21 binaries attributed to a container base
image whose payload could not be downloaded; the catalogue traversal includes 11
binaries from a benchmark fixture that has no scan result. One of those terms
deserves its own mention: an artifact whose payload is absent from disk, and
which the catalogue traversal therefore reports as zero binaries, still
contributes 21 binaries to the published inventory total through a scan summary
written before the payload went missing. That is a stale record feeding a
current figure, and it is the same defect class as the denominator issue above:
two traversals that disagree about what the corpus is.

\subsection{Result 2: D02 on internal symbols is unfalsifiable}
\label{sec:soundness}

Of the 40 findings that reached a decision in the pre guard snapshot, 35 were
cleared \texttt{not\_affected} under D02. All 35 were OpenSSL, spanning three
CVEs. KT-2 required hand verification, and the first statement examined exposed
the defect.

The pipeline claimed that \texttt{OSSL\_CRMF\_ENCRYPTEDVALUE\_decrypt} was absent
from \texttt{libcrypto.so.3} in a Redis container image. The claim was true as
stated. The function is not in the dynamic symbol table. It is also not in the
dynamic symbol table of any OpenSSL build ever produced, because it is internal
and OpenSSL does not export it.

\begin{lstlisting}
$ readelf --dyn-syms -W libcrypto.so.3 | grep -i CRMF
OSSL_CRMF_ENCRYPTEDVALUE_free
OSSL_CRMF_CERTID_new
OSSL_CRMF_CERTTEMPLATE_free
... approximately forty more

$ objdump -d libcrypto.so.3 | grep -ic crmf
597
\end{lstlisting}

The certificate request messaging subsystem is compiled in. There are 597
references to it in executable code. The deallocation function for the exact data
type the vulnerability operates on is exported. Debian does not build with the
relevant feature disabled. The single thing absent is the one function the check
looked for, and its absence was structurally guaranteed before the check ran.

\paragraph{Why this is worse than a wrong answer.} A shared library presents two
distinct surfaces. The export table lists functions other programs may call. The
executable code contains everything the library does, most of which is private.
Our check read the export table and treated absence there as evidence about the
code. Those are different claims, and the great majority of a library's functions
never appear in an export table while remaining present and executable.

The clearance was therefore not merely incorrect. It was \emph{unfalsifiable}:
there exists no binary anywhere for which that check could have returned anything
other than absent. The tool had a rule that always says yes and reported its
output as evidence.

\begin{table*}[t]
\caption{Symbol observability across the three cleared CVEs, with one contrast
case. Only the last row describes a check whose outcome was capable of varying.}
\label{tab:symbols}
\centering
\footnotesize
\setlength{\tabcolsep}{4pt}
\begin{tabular}{@{}llllp{4.4cm}@{}}
\toprule
CVE & Symbol & Library & Exported & Code present \\
\midrule
CVE-2026-42767 & \texttt{OSSL\_CRMF\_ENCRYPTEDVALUE\_decrypt} & \texttt{libcrypto.so.3} & no & 597 CRMF references \\
CVE-2026-45445 & \texttt{aes\_ocb\_cipher} & \texttt{libcrypto.so.3} & 0 & 304 OCB references \\
CVE-2026-14456 & \texttt{port\_default\_packet\_handler} & \texttt{libssl.so.3} & 0 & \emph{Withdrawn, see Section~\ref{sec:version}.} As measured: 691 DTLS references; 13 exported DTLS symbols; 43 DTLS strings. This CVE is QUIC, not DTLS, and does not apply to the branch measured. \\
CVE-2023-0286 & \texttt{GENERAL\_NAME\_cmp} & \texttt{libcrypto.so.3} & \textbf{yes} & contrast case; absence here would be meaningful \\
\bottomrule
\end{tabular}
\end{table*}

Table~\ref{tab:symbols} shows the pattern held across all three cleared CVEs. The
contrast row matters: \texttt{GENERAL\_NAME\_cmp} is part of the public API and
does appear in the export table, so for that CVE the same rule produces genuine
evidence. One symbol in four. An earlier draft stated that the five surviving
\texttt{affected} verdicts all run on this path under D04. That was wrong and is
withdrawn in Section~\ref{sec:version}: they span two CVEs.

Across the corpus the 35 clearances distribute as 22 instances of
CVE-2026-14456, 9 of CVE-2026-42767 and 4 of CVE-2026-45445. Twenty eight of the
35 are unsound on the mechanism above; Section~\ref{sec:version} accounts for
the remaining seven.

\paragraph{Root cause in the manifest.} The manifest entries record these symbols
with \texttt{confidence: definitive} and a \texttt{patch-touched-function}
derivation. That derivation correctly identifies the function the upstream fix
touched. It does not check whether the function is observable in a stripped,
dynamically linked build, nor which shipped shared object hosts it. The note on
one entry states that the function is not exported on the corpus builds and
treats that as supporting evidence. The defect is written into the manifest by
its own author, in plain language, and neither the author nor any check in the
pipeline recognised it.

\subsection{KT-2: fail}

\begin{table}[t]
\caption{KT-2 hand verification, 100 statements. The scored column is the
pre registered result as recorded. The corrected column is a later overlay
described in Section~\ref{sec:version}; the labels themselves were not
rescored.}
\label{tab:kt2}
\centering
\begin{tabular}{@{}lrr@{}}
\toprule
Metric & As scored & Corrected \\
\midrule
Statements labelled & 100 / 100 & 100 / 100 \\
False \texttt{not\_affected} & \textbf{20 / 20} & \textbf{16 / 20} \\
Refusals correct & 80 / 80 & 80 / 80 \\
Overall agreement & 80\% & 84\% \\
\textbf{Verdict} & \textbf{FAIL} & \textbf{FAIL} \\
\bottomrule
\end{tabular}
\end{table}

As scored, all 20 clearances in the labelled set are false and all 80 refusals
agree with the hand label. The pre registered bar was zero false clearances, so
KT-2 fails on the first statement examined and would fail on any of the other
nineteen. The sample distributes as 13 instances of CVE-2026-14456, 6 of
CVE-2026-42767 and 1 of CVE-2026-45445.

Four of those 20 labels were later found to be wrong, in the direction of
harshness: the pipeline's clearance was correct and our hand label was not.
Section~\ref{sec:version} gives the reason. The corrected count is 16 false
clearances, and KT-2 still fails, because the bar was one and not a rate. We
report both columns rather than replacing one with the other, and the labelled
file records the original label alongside a dated amendment rather than
overwriting it.

One reporting point, because this is easy to present misleadingly. Precision and
recall computed per reason code come out at 1.000 across the four labelled
refusal codes. That figure describes refusal rows only. The clearances carry no
reason code and appear in the same table as a row with recall 0.000 and 20 false
negatives. Reporting the 1.000 figures without the adjacent zero row would be an
accurate sentence assembled into a false impression. The two must always be
printed together.

\subsection{Engine response, reported separately}

After the finding, D02 was changed to refuse unless the manifest asserts
\texttt{dynsym\_export\_verified} on every definitive vulnerable symbol. Only
\texttt{GENERAL\_NAME\_cmp} carries that flag. Re running decide over the corpus
converts all 35 clearances to \texttt{symbol\_not\_observable} refusals at S5b,
28 of which were unsound and 7 of which were correct for a reason the pipeline
never computed (Section~\ref{sec:version}).
The post guard pipeline emits zero \texttt{not\_affected} verdicts on this
corpus.

That is the correct behaviour and it is also not a product. A clearance tool that
clears nothing has removed its own reason to exist. It does not restore KT-1
either: 99.60\% of findings still terminate at S1 or S2.

It also, as the next section describes, did not fix the defect underneath the
one it was aimed at.

\subsection{Result 3: version applicability was never computed}
\label{sec:version}

Some days after the soundness finding was written up, we checked each of the
four manifest CVEs against its published advisory. Two of the four do not apply
to the versions on the artifacts where the pipeline decided them.

The clearest case concerns the CVE we had described as a DTLS issue. It is not.
It is an unbounded memory growth defect in a QUIC server component, present only
since the release in which that component was introduced, and its advisory
excludes the entire earlier branch that four of our artifacts ship. On those
artifacts the vulnerable code is genuinely absent, because the subsystem does
not exist in that branch. The pipeline's \texttt{not\_affected} was the correct
verdict, arrived at through reasoning that had nothing to do with the reason it
was correct.

Our own supporting evidence for that CVE had compounded the error. Having
already corrected a wrong library check, we had gone on to count references to
the wrong protocol entirely, on a build with no QUIC server in it, and reported
those counts as proof that the vulnerable code was present.

\paragraph{Three layers, one blind spot.} The scanner reported findings on
artifacts whose observed version its advisory excludes, and on one artifact
reported a package version that both the package metadata and the shipped
library contradict. The engine never compares a finding's version against the
manifest's affected version list before a symbol rule fires; the list is data the
implementation loads for its symbols and never consults for its versions. And
the hand verification, ours, asked whether the symbol evidence supported the
claim without first asking whether the claim was in scope.

The consequences are bounded and worth stating precisely.

\begin{table}[t]
\caption{Reclassification of the 35 pre guard clearances once version
applicability is computed.}
\label{tab:version}
\centering
\begin{tabular}{@{}lr@{}}
\toprule
Category & Count \\
\midrule
CVE applies; symbol internal and unobservable & \textbf{28} \\
CVE does not apply to the observed version & 7 \\
\midrule
Total pre guard clearances & 35 \\
\bottomrule
\end{tabular}
\end{table}

Twenty eight of the 35 clearances remain unsound on the mechanism of
Section~\ref{sec:soundness}. Seven were correct by accident. In the labelled
sample the same split gives 16 false clearances rather than 20. Of the 40 pre
guard decisions, a version gate would have terminated 7 before any symbol rule
ran.

\paragraph{What this does not overturn.} The five surviving \texttt{affected}
verdicts are not version false: one is on firmware running a branch the advisory
does list, and four are on a CVE whose range does include the observed versions.
We had also written that all five rested on a single public API symbol, and that
was wrong; they span two CVEs. One of the five is supported by evidence taken
from a second, differently versioned copy of the library shipped under a separate
prefix on the same image, rather than from the package the finding refers to.
That verdict is correct, and nothing in its evidence bundle demonstrates it.

\paragraph{The shape of the error.} The guard added in response to
Section~\ref{sec:soundness} closed symbol observability. It did not close
version applicability, library home binding, or which copy of a library an
observation came from. A correction aimed one layer too shallow leaves the
record looking repaired while the deeper defect is untouched, and we would not
have found this one either had we not gone back to the advisories after the
write up was already drafted.

\subsection{D01: no corpus instance}

D01 is untouched by the soundness finding, because component presence is
observable on stripped dynamic binaries through identity symbols and library
dependencies. It is also untested, for a reason worth reporting.

The corpus produced zero D01 rows. Of 55 scanned directories, 22 carry a Grype
match on a package whose name or PURL names OpenSSL. A path agnostic search
locates an OpenSSL shared object on all 22, and on none is one absent. There was
no natural instance in which the scanner reported a component the artifact did
not contain.

That absence is informative: scanner component attribution was correct on
presence in every case we could test. The open question throughout was whether
the vulnerable \emph{code} was present, which is exactly what D02 could not
answer. Only a synthetic probe on a BusyBox image exercises D01, and the honest
claim is that D01 did not fail on the single artifact tested and remains untested
at scale.

Note the three way count this depends on: 22 directories with a Grype package
match, 19 that ever bound \texttt{component=openssl} in the pipeline, and 55
scanned. The three RPM images appear in the 22 but not the 19, which is Result 1
visible in a second view of the same data.

\section{A secondary result: one scanner, two input paths}
\label{sec:interop}

While assembling the corpus we encountered an unrelated defect with more
immediate practical consequence than anything above. It concerns Trivy 0.72.0
reading SBOMs produced by Syft 1.51.0, two widely deployed open source tools.
Both findings below were reported upstream with complete reproductions on public
images before this paper was prepared, as False Detection discussions numbered
11139 and 11140 on the Trivy repository, and neither was withheld pending
publication.

\subsection{CycloneDX: partial loss, in both directions}

Scanning a container image tar directly yields 409 unique CVE IDs. Generating a
CycloneDX SBOM of the same tar, then scanning that SBOM with the same Trivy
binary against the same database on the same day, yields 137. No error is
emitted, no warning appears, and the SBOM is schema valid.

Debian files security advisories against source packages rather than the binary
packages that install. Syft records the source package twice, as a
namespaced property and as a PURL qualifier. Trivy's SBOM ingest path reads
neither; it looks for its own namespaced property and, in its absence, silently
sets the source name to the binary name. Every Debian advisory match then fails.

The disagreement is bidirectional, and this matters for how the result should be
stated. The image result contains 278 IDs absent from the SBOM result. The SBOM
result contains 6 IDs absent from the image result. The numeric difference
between the two counts is 272, which is 278 minus 6. Copying the value Syft
already wrote into the property Trivy reads produces 415, which
is exactly the union of the two sets. So the SBOM path is not simply a lossy view
of the direct path, and neither result is a superset of the other. Any framing
that treats the direct scan as ground truth is overstating what the data shows.

Across 22 public images, 14 are affected. The most extreme moves from 335 to
4,007 after correction. On three Debian images the count of packages whose source
name differs from the binary name is zero out of 149, zero out of 112 and zero
out of 88 through the SBOM path, against 114, 85 and 62 through the direct path.
Zero on every image is not a sampling artifact. The eight unaffected images are
APK, RPM or Go binary based, and that split along Debian's advisory model is the
evidence for the mechanism rather than the list itself. Four of those eight
report zero on both paths and are therefore uninformative on their own; the
control that carries weight is Grype, which matches its own direct
scan on every image and both formats we tested.

Two clarifications. The mechanism was already documented by a maintainer for a
different distribution family in 2024, and we credit that rather than claiming
discovery; our contribution is the effect size, the control group and the
property copy intervention. And we make no claim about which count is correct for
triage purposes, since many recovered findings carry a distribution status that
de prioritises them. The claim is that two paths through one tool disagree
without saying so, and the information required to reconcile them is already in
the file.

\subsection{SPDX: the OS is never detected}

Syft asked for SPDX rather than CycloneDX, on the same tar,
produces a document carrying 151 packages with valid PURLs. Trivy retains
one and reports zero vulnerabilities. The surviving package is a genuine Java
component rather than a misread system package.

The debug output identifies the mechanism, and it is not PURL parsing. On the
SPDX path Trivy logs that it detected no operating system, warns that the
family is unsupported, and never runs OS package detection at all. On the
CycloneDX path, from the same Syft run on the same tar, it detects the
distribution and version correctly and proceeds with 149 packages. A second
Debian image shows the same shape: 89 packages with PURLs in the SPDX document,
zero packages and zero vulnerabilities out, against 87 unique vulnerabilities
from the direct scan.

This reproduces across Syft versions sixteen months apart. Grype
reading the same SPDX documents retains the packages and matches its own direct
scan, so the information is present and consumable.

The compliance consequence is direct, and it differs in kind from the partial
under count above. An organisation that standardises on SPDX for its regulatory
deliverables may receive a clean bill of health from a tool that skipped its
entire system inventory, from an artifact that is complete, valid and correct. A
zero is indistinguishable from a genuinely clean image from outside.

\section{Discussion}

\subsection{Two independent reasons to stop}

KT-1 asks whether the pipeline can decide at scale, and the answer is that we do
not know, because it stops at identity. KT-2 asks whether its decisions are
sound, and the answer is no, on a mechanism we can name. Either alone would
justify stopping under the pre registered rule. Together they say something more
specific than a general negative.

\subsection{What symbol absence can and cannot support}

The generalisable claim is narrow and, we think, useful.
\texttt{Vulnerable\_code\_not\_present} justified by symbol table absence is sound
only for symbols that would appear in the export map of a reference build. For
internal functions, and internal functions are the majority, absence is
guaranteed before the check runs and the tool will confidently clear everything
it is pointed at.

Anyone automating this needs a manifest field asserting observability, verified
against a reference build, and needs to refuse rather than clear when it is
absent. They also need to record which library hosts each symbol. Our engine
searched every inventoried binary's export table indiscriminately, which allowed
a DTLS symbol to be checked against the wrong library, since DTLS does not live
where the rest of the cryptographic primitives do. A check against the wrong file
is not a negative result for that CVE.

Beyond that, internal functions require disassembly level evidence of absence,
which is materially harder than reading a symbol table and is not addressed here.

\subsection{Identity, not evidence, is the binding constraint}

The field's attention, including ours, has been on the hard and interesting
problem of proving code absent from a binary. Our measurement says that is not
where the pipeline breaks. It breaks on establishing that a scanner's package
reference and an advisory's component are the same thing, that a source package
name in one ecosystem corresponds to an upstream project, and that a version
string maps to a version range. Until that is solved the interesting question
does not arise, and the per ecosystem non uniformity in
Section~\ref{sec:results} suggests it must be solved more than once.

\subsection{Two errors we made finding this}

Both are in the record deliberately, because a tool whose premise is that it only
claims what it can prove should be candid about how easy it is to prove the wrong
thing.

We checked the DTLS symbol against the wrong library first, found zero
references, and nearly wrote it up as a correct clearance. That would have been a
wrong conclusion reached by a wrong method that happened to agree with the answer
we wanted. Separately, hunting for artifacts where OpenSSL was genuinely absent,
we searched a Debian multiarch path glob and got three false candidates including
a firmware image, before a path agnostic search cleared them.

The shared defect class is assuming file layout instead of asking the filesystem.
It cost two near misses in one afternoon on a project whose entire premise is
deterministic evidence, and Section~\ref{sec:denominator} shows a third instance
of the same class at the level of the corpus itself.

\subsection{The uncomfortable part}

The tool did not catch its own error. Nothing in its design was capable of
noticing, because the failure was not a bug in a rule but a category error in
what the rule's input could mean. Detection required a person with
\texttt{readelf} and an afternoon of suspicion.

That is not a comfortable conclusion for anyone automating evidence production
for regulatory filings, and we do not have a general answer to it. The narrow
answer is that hand verification against ground truth is not optional, and that a
pre registered soundness bar of zero is what forced us to look.

\section{Limitations}
\label{sec:limits}

\begin{enumerate}
\item \textbf{Single implementation.} One engine, not independently
reimplemented.
\item \textbf{Single scanner for the corpus.} Findings come from Grype 0.115.0
only.
\item \textbf{OpenSSL only soundness path.} All 35 clearances and all labelled
false clearances are OpenSSL CVEs. Nothing here establishes unsoundness for
other components or for D01 and D04.
\item \textbf{Version strings and backports.} Per artifact versions come from
the string in the shipped library, corroborated by package metadata. A
distribution that backports a fix without changing the upstream version string
would defeat that oracle, and Debian routinely does exactly this. Our
applicability determinations are therefore sound for the case where a version is
outside an advisory's range entirely, which is the case we rely on, and weaker
for the case where a version sits inside a range but may carry a backported fix.
\item \textbf{No version gate.} The engine never computes version
applicability (Section~\ref{sec:version}). Every clearance and every
\texttt{affected} verdict in this paper was produced without it.
\item \textbf{Which copy was measured.} The engine binds a symbol observation to
any inventoried binary rather than to the library the finding names. One of the
five \texttt{affected} verdicts is supported only by observations on a second
copy of the library at a different version.
\item \textbf{D01 untested at scale.} Zero natural corpus instances; one
synthetic probe.
\item \textbf{Frozen KT-2 snapshot.} Scores the pre guard pipeline by design
(Section~\ref{sec:frozen}); post guard figures reported separately.
\item \textbf{Denominator composition.} The KT-1 denominator includes 253
findings from two directories that are not product artifacts
(Section~\ref{sec:denominator}).
\item \textbf{Stale record in a current figure.} The published inventory total
includes 21 binaries from an artifact whose payload is absent from disk
(Section~\ref{sec:denominator}).
\item \textbf{Manifest depth.} 25 seed components, of which 7 are provenance
verified. Most scanner CVE and component pairs have no manifest entry at all,
which is itself part of the S2 count.
\item \textbf{Identity alias immaturity.} Ten alias rows, nine probable, one hand
verified definitive. The identity layer is a seed, not a mature mapping, and
Result 1 is in part a measurement of that immaturity rather than of the ecosystem
alone.
\item \textbf{Layout assumptions.} Engine and analysis scripts embedded Debian
multiarch path priors, corrected during this work.
\item \textbf{Firmware stratum size.} 12 artifacts, one decided finding.
\item \textbf{Static and musl release binaries.} Eight artifacts, zero decided.
\item \textbf{Substitute artifacts.} Seven container stand ins for firmware
targets that could not be obtained, retained with honest class labels and
excluded from firmware stratum claims.
\item \textbf{Three instrument passes.} Three consecutive not evaluable readings;
the pre registered rule treats this as terminal rather than as an invitation to a
fourth.
\item \textbf{Interoperability tooling provenance.} Binaries used for
Section~\ref{sec:interop} live outside the repository; hash verification against
official checksums is the mitigation. Those figures were computed from cached
SBOMs and re verified against a current vulnerability database at time of
writing.
\item \textbf{No refitting.} No threshold moved and no rule was refitted to
observed data. Changes needed for a second version are recorded as future work.
\end{enumerate}

\section{Ethics and responsible disclosure}
\label{sec:ethics}

\paragraph{Disclosure.} Both interoperability defects in
Section~\ref{sec:interop} were reported to the maintainers of the affected
project, with complete reproductions on public images, before this paper was
prepared for publication. Each report credits the prior art we found for the
underlying mechanism and states plainly which parts of our contribution are new.
Where an arithmetic error was found in our own filed report after posting, we
published a correction on the same thread rather than editing the original, so
the record shows the correction.

We characterise both as interoperability gaps between namespaced metadata
conventions rather than as defects in either tool, and we take no position on
which of the two disagreeing counts is correct for any given triage purpose.
Neither finding enables an attack: both cause a scanner to report fewer findings
than another path through the same tool, which is a reporting fidelity problem
for the operator, not a capability granted to an adversary. We did not withhold
either report pending publication.

\paragraph{No new vulnerabilities.} This work reports no previously unknown
software vulnerability. Every CVE referenced was already public at the time of
measurement. The soundness finding in Section~\ref{sec:soundness} is a defect in
our own tool's reasoning, disclosed here in full.

\paragraph{Scanning third party artifacts.} All artifacts were obtained from
publicly available distribution channels: vendor firmware download pages, project
release pages, and public container registries. Analysis was limited to static
inspection of binaries already present on disk. No device was operated, no
network service was probed, no access control was circumvented, and no
proprietary artifact is redistributed by us.

\paragraph{Human subjects.} None. This work involves no human participants and no
personal data, and required no institutional ethics review.

\section{Data availability and licensing}
\label{sec:data}

\paragraph{What we publish.} The corpus catalogue, recording for every artifact
its identifier, class, source URL, fetch timestamp and SHA-256; the complete
decision record for every finding; the evidence bundles for decided findings; the
hand labelled ground truth; the computed metrics; and the figure provenance
audit. Code and re derivation scripts are published under Apache 2.0; the data
under Creative Commons Attribution 4.0.

\paragraph{What we do not publish, and why.} We do not redistribute any scanned
artifact. The corpus contains proprietary vendor firmware, container images and
vendor supplied toolchains whose licences do not permit redistribution by a third
party. The catalogue instead records a source URL and a cryptographic hash for
each, so a reader can obtain the same bytes from the original publisher and
verify they match what we measured. Where an artifact could not be obtained at
all, the catalogue records it as absent rather than substituting something else;
where a stand in was used it carries an explicit class label.

This means reproduction depends on upstream availability. Vendor download pages
move and release archives are withdrawn. The recorded hashes let a reader confirm
that whatever they retrieve is what we analysed, but we cannot guarantee
retrieval itself. We regard this as the honest trade against redistributing
material we have no right to redistribute.

\paragraph{Trademarks.} Product, project and company names referenced here are
the trademarks of their respective owners and are used descriptively to identify
the software measured. Their use implies no affiliation with or endorsement by
those owners.

\paragraph{Not legal advice.} This paper describes a regulatory context to
motivate a technical problem. Nothing in it is legal advice, and nothing in it
establishes whether any particular product, statement or artifact satisfies any
obligation under the Cyber Resilience Act or any other regime. Responsibility for
statements in a technical file rests with the manufacturer.

\section{Conclusion}

We built a binary grounded VEX clearance pipeline, pre registered two kill tests,
and both resolved against it.

The coverage test is not evaluable rather than failed, because 99.60\% of
attributed findings terminated at identity resolution or manifest lookup before
any binary was opened, and zero terminated at a symbol table stage despite 83\%
of corpus binaries being stripped. The soundness test fails absolutely: of the 35
clearances the pipeline emitted, 28 rested on absence of internal functions from
an export map where absence was guaranteed a priori, and the remaining 7 were
correct only because the CVE did not apply to the shipped version, which the
pipeline never checked.

The useful residue is not a tool. It is a stage attributed map of where this
approach dies on real shipped artifacts, a named failure mode that anyone
building symbol based clearance should know about before shipping, and a hand
labelled ground truth that did not previously exist publicly.

The corpus catalogue with hashes, every decision the pipeline made, the evidence
bundles and the labelled statements are published. Every figure in this paper
regenerates from a single script, and the discrepancies between the two
traversals of Section~\ref{sec:traversals} are published as an audit record
rather than reconciled away. The material required to prove us wrong is included
deliberately.


\begin{acks}
We thank the maintainers of the open source scanning tools measured here. Both
defects reported in Section~\ref{sec:interop} were filed upstream with
reproductions before publication.
\end{acks}

\bibliographystyle{ACM-Reference-Format}
\begin{thebibliography}{9}

\bibitem{churakova}
Y. Churakova, M. Ekstedt, and L. Schmid.
\newblock Vexed by VEX tools.
\newblock arXiv:2503.14388.

\bibitem{rasheed}
S. Rasheed, M. McPhee, L. Patterson, S. MacDonell, and J. Dietrich.
\newblock Hidden dependencies and component variants in SBOM based software
composition analysis.
\newblock arXiv:2604.21278.

\bibitem{cofano}
D. Cofano et al.
\newblock On SBOM precision and vulnerability loss in the Python ecosystem.
\newblock CITATION TO BE COMPLETED FROM THE PRIMARY SOURCE.

\bibitem{benedetti}
G. Benedetti et al.
\newblock On the identifiability of vulnerabilities from generated SBOMs.
\newblock CITATION TO BE COMPLETED FROM THE PRIMARY SOURCE.

\bibitem{dann}
A. Dann et al.
\newblock Evaluating software composition analysis tools.
\newblock CITATION TO BE COMPLETED FROM THE PRIMARY SOURCE.

\end{thebibliography}

\end{document}
